% TOP_PROBLEM: {"id": 8700014, "title": "Conjecture 2.6", "category_id": 3, "category": {"id": 3, "name": "graph_theory", "display_name": "Graph Theory", "description": "Problems involving graphs, networks, and their properties.", "slug": "graph-theory", "order_index": 3, "created_at": "2026-07-31T15:26:25.671Z"}, "status": "open", "problem_number": "AMR-086-0014", "set_id": 13}
% FIRST_POSTED: 2026-10-01
% ENTRY_KIND: statement_only
% UnsolvedMath ID 8700014; source catalogue code AMR-086-0014
% !TeX program = lualatex
% Definition and sources only; this file is not an LLM solution attempt.
\documentclass[11pt,a4paper]{article}
\usepackage[margin=25mm]{geometry}
\usepackage{fontspec}
\setmainfont{lmroman10-regular.otf}[BoldFont=lmroman10-bold.otf,
 ItalicFont=lmroman10-italic.otf,BoldItalicFont=lmroman10-bolditalic.otf]
\usepackage{amsmath,amssymb,mathtools}
\usepackage{unicode-math}
\setmathfont{latinmodern-math.otf}
\AtBeginDocument{\let\setminus\smallsetminus}
\usepackage{xurl,hyperref}
\hypersetup{colorlinks=true,urlcolor=blue}
\setcounter{secnumdepth}{0}
\setlength{\parindent}{0pt}
\setlength{\parskip}{5pt}
\setlength{\emergencystretch}{3em}
\begin{document}
\section*{Conjecture 2.6}\label{8700014}
{\small UnsolvedMath ID 8700014 · AMR-086-0014 · Graph Theory · AMR Open Problem Lists}
\subsection{Definitions and mathematical statement}
Let $x,y,p,q$ be arbitrary positive integers, and suppose $x^p\ne y^q$. Put $M=\max\{x^p,y^q\}$, and interpret powers of $M$ with real exponent in their usual positive-real sense. The proposed inequality is
\[
|x^p-y^q|\ge C(\varepsilon)
M^{\,1-1/p-1/q-\varepsilon}.
\]
The question is whether, for every real $\varepsilon>0$, a constant $C(\varepsilon)>0$ exists for which this inequality holds simultaneously for every permitted quadruple $(x,y,p,q)$.

The order of quantifiers matters: $C(\varepsilon)$ may depend on the chosen error exponent, but not on either base, either exponent, or their greatest common divisor. The bases need not be coprime, and the exponents may be equal. The statement includes $p=1$ or $q=1$; in those elementary cases the right-hand power is nonpositive, but retaining them makes the domain exactly the positive-integer domain of the source.

Equality of the two powers is excluded, since it would contradict any positive lower bound. When they differ, the left side is a positive integer, which gives only the basic bound one. The conjecture predicts substantially growing gaps in the regimes where $1-1/p-1/q-\varepsilon>0$.

No formula or effective computation for the constant is demanded. The target is its existence with the stated uniformity, rather than separate inequalities whose constants vary with $p$ and $q$.
\subsection{Short English statement}
Are unequal positive integer powers separated by the stated power-sized gap, with a constant depending only on the error exponent?
\subsection{Sources}
\begin{itemize}
\item[S1] ULAM.AI and the UnsolvedMath Contributors, supplied problems.json, record 8700014 (AMR-086-0014), AMR Open Problem Lists. Catalogue identity and source transcription; CC BY 4.0.
\url{https://huggingface.co/datasets/ulamai/UnsolvedMath}.
\item[S2] Waldschmidt - Open Diophantine Problems (2004), Conjecture 2.6. Original source indexed by AMR-086-0014.
\url{https://arxiv.org/abs/math/0312440}.
\end{itemize}
\end{document}
